\section{Coefficient bounds and the remaining complexity barrier}
\label{sec:complexity}

\subsection{An explicit bound for an actual completed suffix}

The completion theorem is useful only with a proved size bound. Let $n\ge1$
be the number of suffix crossings. A completed suffix is obtained by pairing
its frontier labels, with no additional crossing-free circles supplied from
outside. Circles already removed by delooping in the prefix are represented by
the relative complex and its shifts; they must not be silently added to this
closed diagram.

\begin{lemma}[Crossings bound the number of link components]
\label{lem:components-crossings}
For a spherical $n$-crossing diagram in which every link component meets a
crossing, the number $c$ of link components satisfies $c\le n$.
\end{lemma}
\begin{proof}
For component $i$, let $s_i$ be its number of self-crossings and $u_i$ the number
of mixed-crossing incidences on it. If $s_i>0$, then $s_i+u_i/2\ge1$.
If $s_i=0$, the component is a simple closed curve on the sphere. Its number of
intersections with every other closed curve is even, by mod-two intersection
on the sphere. Because it meets a crossing, $u_i$ is positive and at least two.
Thus $s_i+u_i/2\ge1$ in this case as well. Summing over components counts
each self-crossing once and each mixed crossing twice with coefficient one
half, so $c\le\sum_i(s_i+u_i/2)=n$.
\end{proof}

Every link component of our completed suffix meets a retained crossing:
each frontier endpoint belongs to a retained crossing, and the matching has
no closed component of its own. Hence the lemma applies even if the projection
is disconnected.

\begin{proposition}[Per-object and per-component coordinate bounds]
\label{prop:coordinate-bounds}
Let $\sigma\in\Z^4$ be the marked four-residue Euler vector of a completed
$n$-crossing suffix object, and put $a=J(1)$. Then
\begin{align}
 |a|&\le 2^{n-1}, & |d|&\le 2^{n-1},\\
 \norm{\sigma}&\le2^{n-1}, &
 |\sigma_r|&\le\left\lfloor\frac{|a|+2^{n-1}}2\right\rfloor.
 \label{eq:object-bounds}
\end{align}
For a whole differential component containing $s_a$ represented objects,
with object Euler values $a_o$, a valid bound is
\begin{equation}
 C_a=\sum_{o\in\mathcal C_a}
       \left\lfloor\frac{|a_o|+2^{n-1}}2\right\rfloor
       \le s_a2^{n-1}.
 \label{eq:coordinate-bound}
\end{equation}
\end{proposition}
\begin{proof}
The value-at-one formula gives $|a|=2^{c-1}\le2^{n-1}$. For a connected
projection, the signed Tait cofactor is a sum of products of $\pm1$ over
spanning trees. If its Tait graph has $v$ vertices, every tree chooses exactly
$v-1$ of the $n$ crossing edges, so its absolute value is at most
$\binom{n}{v-1}$. For $n\ge1$, each binomial coefficient is at most
$2^{n-1}$: it is a term in either the even or the odd binomial sum, and both
sums equal $2^{n-1}$. Loops and parallel edges do not invalidate this argument.
For a disconnected projection, the value at $i$ is zero, so the same bound
holds. Multiplication by a Gaussian unit does not change $|d|$.

The two nonzero possible coordinates are $(a+d)/2$ and $(a-d)/2$; therefore
their total absolute value is $\max(|a|,|d|)$, and each absolute value has
the stated integer bound. Homological signs and quantum shifts only negate
or permute coordinates. Summing over a whole component and using the triangle
inequality gives \cref{eq:coordinate-bound}.
\end{proof}

This deliberately elementary bound is inexpensive to compute and is valid for
signed, nonalternating diagrams. A smaller certified bound can only improve
the stopping criterion. The implementation uses the first expression in
\cref{eq:coordinate-bound}, because the exact $a_o$ values have already been
computed by the Euler prefilter.

\subsection{How many bits of modulus suffice?}

Let $s_{\max}=\max_a s_a$. At threshold one, a product with
\[
                   M>s_{\max}2^{n-1}
\]
is sufficient. Thus the sufficient threshold has $n+\log_2 s_{\max}+O(1)$ bits.
A small-prime product can overshoot that target by the bit length of its last
prime; this does not affect the asymptotic bound for the Boolean comparison. No exponential-size
integer \emph{bit length} is hidden in this step, although $s_{\max}$ itself
may count exponentially many represented objects.

For a supply of distinct primes each of at least $w$ bits, the number needed
is at most
\[
 O\!\left(1+\frac{n+\log(s_{\max}+1)}{w}\right),
\]
provided enough such primes are supplied and validated. This is a statement
about the arithmetic interface, not a claim that a fixed machine-word palette
works for arbitrarily large inputs. An asymptotic implementation can enlarge
the prime range and include the cost of finding and certifying the primes.
For completeness, let $L\ge1$ be the required modulus bit length. Sieve
odd primes up to a cutoff $X$, doubling $X$ until their product exceeds the
certified bound. The classical Chebyshev estimate
$\vartheta(X)=\sum_{p\le X}\log p=\Theta(X)$ \cite{chebyshevnotes}
shows that this stops with $X=O(L)$. Sieving, exact product construction,
and primality certification by the sieve therefore take polynomial time and
space in $L$, and each selected prime has $O(\log(L+1))$ bits. This supplies
a deterministic asymptotic extension; it is not implemented in the fixed
palette interface. A user can also supply a sufficiently large certified
finite palette for a specified input-size range.

The current code accepts distinct odd primes below $2^{31}$, and defaults to
$(65521)$. It tests the explicit completion inequality after every full prime
pass. If the supplied palette is exhausted before that inequality holds, the
record says that threshold completion was not certified. Scanning then
continues. The code neither invents missing modulus bits nor interprets an
inconclusive residue as an unknot certificate.

\subsection{A cost account for the implemented observer}

Let $N$ be the original number of crossings. For observed suffix $s$, write
$n_s$ for its crossing count, $b_s$ for its frontier size, $v_s$ for the number
of common black-fragment vertices, $\ell_s$ for the number of primes processed,
$K_s$ for the number of distinct boundary matchings queried, and $D_{s,p}$
for the number of distinct black terminal partitions queried at prime $p$.
Let $S_s$ be the total number of represented objects traversed in the observed
whole components. These are measured workload parameters, not promised small
functions of $N$.

The main costs are as follows. Ordinary dictionary and disjoint-set operations
are stated in the usual unit-cost model; field and large-integer costs are
separated explicitly.

\begin{center}
\begin{tabularx}{\textwidth}{@{}P{0.29\textwidth}X@{}}
\toprule
Operation & Cost and limitation\\
\midrule
Source coloring & Once, near-linear in $N$, including spherical validation.\\
Suffix geometry & $O(N+n_s\log(n_s+1)+v_s^2)$ elementary work is a safe bound.
The current constructor revalidates the whole crossing order at each stage;
it is not a suffix-only $O(n_s)$ routine.\\
Boundary acceptance and counts & $O(b_s\log(b_s+1))$ per new matching is a safe
bound for normalization and boundary disjoint-set summaries.\\
Response preparation & $O(v_s^3)$ field operations and $O(v_s^2)$ temporary
field elements per prime.\\
New determinant query & $O(1+b_s^3)$ field operations, using a matrix of size
at most $b_s-2$ when $b_s>0$; immediate zero when radical dimension is too large.\\
Whole-component aggregation & $O(\ell_sS_s)$ four-coordinate visits, in addition
to grading recovery and exact Euler accounting. No norm is applied to a proper
subset of one component.\\
CRT and norm minimization & Four-coordinate integer arithmetic per component
and prime, on integers whose bit length includes $\log M$ and the exact sums.\\
\bottomrule
\end{tabularx}
\end{center}

In particular, the determinant field-operation term is bounded by
\begin{equation}
 O\!\left(\sum_s\left[
       \ell_s v_s^3+\sum_{p\ {
m processed}}D_{s,p}(1+b_s^3)
                      \right]\right).
 \label{eq:arithmetic-ledger}
\end{equation}
This term must be added to the scan itself, geometric validation, grading
recovery, coefficient accounting, cache operations, and CRT bit costs.
If prime bit lengths and the total modulus bit length are polynomial in $N$
and $\log(S_s+1)$, each displayed elementary arithmetic operation has polynomial
bit complexity. This suffices for the conditional asymptotic conclusions below;
we do not equate a Python big-integer operation with a constant-time bit operation.

Active nonzero response data use $O(\ell_s(1+b_s)^2)$ field elements in addition
to the common dense graph. The implementation also retains per-stage matching
and cofactor caches, and completed value caches whose keys include earlier
stages. Its \emph{total} memory is therefore not bounded by $O(b_s^2)$.
Memory accounting must include at least the recorded matching keys, four-vectors,
Euler values, and the scanner's own objects and morphisms.

\subsection{What a width bound would and would not buy}

There are at most
\[
                 (b-1)!!\le b^{b/2}=2^{O(b\log b)}
\]
perfect matchings on a labeled even frontier. Black terminal partitions are
also bounded by the Bell number on $b/2$ elements, hence by
$2^{O(b\log b)}$. These bounds require no single-disk assumption. If all
completions are known to be noncrossing on one fixed disk boundary, a Catalan
bound may be substituted; the present general suffix representation does not
establish that additional hypothesis.

Consequently a polylogarithmic frontier bounds the \emph{number of possible
oracle keys} quasi-polynomially. For example,
$b=O((\log N)^2/\log\log N)$ gives at most $2^{O((\log N)^2)}$ matching keys.
The stronger restriction $b=O(\log N)$ gives
$2^{O(\log N\log\log N)}$ keys by this elementary estimate. These are oracle
table bounds, not complete recognition bounds.

Many algebraic generators may have the same boundary matching and different
differentials, gradings, or continuation behavior. Determinant caching removes
repeated observation cost for that matching; it does not identify those
algebraic generators. Moreover, crossing order optimization has no proved
polylogarithmic-width guarantee on arbitrary inputs in the maintained program.
Planarity by itself is not such a guarantee.

\begin{proposition}[A precise conditional route to quasi-polynomial time]
\label{prop:conditional}
Consider a family of validated diagrams with $N$ crossings. Suppose there is
an algorithm which, through a decisive marked obstruction or final rank
calculation, constructs and updates the required relative complexes in
$2^{(\log N)^{O(1)}}$ bit operations and space, including every represented
object and morphism. Suppose also that the number of stages and oracle calls
is quasi-polynomial and that exact sums and multiplicities have polynomial or
quasi-polynomially bounded bit length. Then adding the deterministic modular
observer with sufficient certified modulus preserves a quasi-polynomial bound.
\end{proposition}
\begin{proof}
A suffix has $O(N)$ darts and at most $O(N)$ common graph vertices. Preparation
and each determinant therefore have polynomial dimension. The coefficient
bound requires $O(N+\log(S_s+1))$ modulus bits; the hypothesis makes this
polynomial in $N$ for a quasi-polynomial number of explicit objects.
The sieve construction above supplies enough primes with polynomial bit
length in polynomial time in the required modulus bit length. Every observer
operation has polynomial cost in its relevant bit lengths, hence at most
quasi-polynomial cost in $N$, apart from its already counted traversal of
represented objects. A quasi-polynomial number of such operations and the
assumed representation work remain quasi-polynomial.
\end{proof}

The main hypothesis is exactly the unresolved obligation. The proposition is
an interface theorem explaining where a future representation result would
attach, not a proof that this hypothesis holds. Even a constant-size determinant
cannot pay for exponentially many distinct differential generators. Conversely,
a structural family formula may succeed despite an exponentially large explicit
basis, as the three-braid literature demonstrates \cite{kelomaki}.

There is a second independent obstruction: the four-residue bound can remain
one for every prefix examined. Khovanov unknot detection concerns the full
homology, whereas these Euler values retain only a small part of it. A universal
bound on the first decisive prefix would need a new theorem about continuation
behavior. None of the modular arithmetic results supplies that theorem.
